\section{Extension: energy loss in a non-rigid body}
\label{sec:extension}
So far the body has been almost perfectly rigid, and its energy has been conserved.
A real body is not perfectly rigid: it flexes slightly as it spins, and internal friction turns some of the energy of this wobbling into heat.
As an extension of the project we study what this energy loss does to the rotation.

\paragraph{The major axis rule}
No external torque acts, so $\vb{L}$ is still exactly conserved and only the energy $E$ decreases.
For a fixed $L$, the energy of rotation about a principal axis is $E=L^2/2I$, which is smallest for the largest $I$.
A body that keeps losing energy at fixed $L$ must therefore end up spinning about its major axis, whatever axis it started on.
This means that for a real, slightly flexible body only the major axis is truly stable: even spin about the minor axis, which is stable for a rigid body (\cref{sec:stability}), slowly drifts away.
This is what happened to Explorer 1 in 1958, and it is why the Earth, which spins about its major axis, keeps a steady rotation.

\paragraph{Model}
The spring model can show this directly, because its body is not perfectly rigid.
We give every spring a damping coefficient $c$, the second term in \cref{lst:spring}, which resists the two masses moving apart or together.
When the body turns rigidly the springs keep their length and the damping does nothing; only the wobbling is damped.
Like the spring force, the damping force is central and equal and opposite, so it removes energy but conserves momentum and angular momentum exactly.
This gives a useful check of the code: $\abs{\vb{L}}$ must stay constant while $E$ falls.

\begin{sloppypar} % units and symbols here would otherwise stick out of the column
We use soft springs, $k=\SI{1}{N/m}$, and $c=\SI{0.7}{N.s/m}$, so that the body wobbles noticeably and settles within a few hundred seconds, instead of the far longer time a nearly rigid body would take.
The soft springs also allow a larger time step, $\Delta t=\SI{5}{ms}$.
The octahedron starts spinning at \SI{1}{rad/s} about $z$, its minor axis, tilted slightly towards $x$, $\vb*{\omega}_0=(0.1,0,1)\,\si{rad/s}$, and is followed for \SI{1000}{s} with both engines.

\paragraph{Numerical details}
With a velocity-dependent force, velocity Verlet is no longer exactly symplectic, which is expected since the energy is meant to fall.
The damping force computed at the end of a step uses the velocity available at that point, the half-step velocity $\vb{v}^{n+1/2}$ (\cref{lst:verlet}).
This adds a small error, of relative size about $c\Delta t/m\approx0.004$ per step.
\end{sloppypar}

\paragraph{Expected outcome}
The initial angular momentum is $\vb{L}=(I_x\omega_x,0,I_z\omega_z)=(2.6,0,10)\,\si{kg.m^2/s}$, so $L\approx10.3\,\si{kg.m^2/s}$ and $E_0\approx\SI{5.1}{J}$.
If the body ends up spinning about $x$, it should turn at $\Omega_x=L/I_x\approx\SI{0.40}{rad/s}$ with $E=L^2/2I_x\approx\SI{2.05}{J}$.
These values are only approximate, because the soft springs let the body stretch and change its moments of inertia while it spins.
The run is shown in \cref{fig:vis-major} and analysed in \cref{sec:res-major}.
